\phantomsection\addcontentsline{toc}{section}{Source index}\section*{Source index}\label{source-notes}

\begin{studycolumns}
\vspace{0pt}
Lecture citations (L1--L18), tutorial numbers and section references throughout identify the evidence. All source PDFs are in the adjacent \code{lectures/}, \code{tutorial/} and \code{pyps/} folders. The editable study guide is \code{../SC2006\_Study\_Guide.md}.
\begin{itemize}
\item \textbf{Chapters 1--6:} Introduction; Requirements Elicitation 1--3; Conceptual Models; Dynamic Models; Writing SRS; UMLSpecNotes.
\item \textbf{Chapters 7--8:} Software Processes -- Basics; Agile Software Development -- Basics, including its image slides.
\item \textbf{Chapters 9--11:} System Design (both decks), Object Design (both decks), Design Patterns, MVC (both versions). The file named \code{Lecture\_13\_System\_Design\_2026.pdf} opens as Object Design.
\item \textbf{Chapters 12--13:} Test Fundamentals/Testing Introduction; Equivalence Class Testing; Boundary Value Testing (2026 and other supplied versions).
\item \textbf{Chapters 14--15:} Tutorial Questions and the scanned GES supplement. Tutorial-only examples remain in Chapter 15; shared answers appear once in their topic chapter.
\item \textbf{Chapter 16:} Understanding AI Engineering; AI Engineering as a practice -- NTU; Agentic Software Engineering; Thursday System Design extension. These guest topics have no separately confirmed assessment status in the supplied files.
\end{itemize}
\textbf{Source conflicts:} use hollow-triangle UML generalization; keep the course's historical Scrum wording distinct from current-standard claims; distinguish graph bases from feasible traces; use the problem's stated mortgage upper limit (8000 versus 80000); \code{800001} in the separate summary is a typo for \code{80001}. Full BV triplets and reduced course selections are labelled separately. Local GES and paper-specific assumptions remain beside the relevant answer.
\end{studycolumns}
